\begin{proposition}[Two failures of interchange in a proposed symbol proof]
\label{prop:B-limit-diagnostics}
\status{Proved counterexamples to two limit arguments; the B1 gate is unchanged.}
The following two implications, which might otherwise be used to attack
Proposition~\ref{prop:B-diagonal-Chow-gate}, are false.
\begin{enumerate}
\item Put \(k=\overline{\F}_2\), \(L=k(t)\), and \(\ell=3\).
The natural map
\[
 \varinjlim_{U\subset\mathbf P^1_k}
 H^1_{\mathrm{cont}}(U,\Z_3(1))
 \longrightarrow H^1_{\mathrm{cont}}(\Spec L,\Z_3(1)),
\]
where \(U\) runs over dense opens, is injective and not surjective.
It remains nonsurjective after tensoring with \(\Q_3\), and an
explicit omitted class is invariant under \(\Gal(k/\F_2)\).
\item There is a torsion-free \(\Z\)-module \(M\) with an automorphism
\(F\), such that every element of \(M\otimes\Q\) has finite
\(F\)-orbit, but the rationalized \(3\)-adic completion has a nonzero
\(4\)-eigenvector.  Thus local finiteness before completion does not
exclude a non-root-of-unity eigenvalue after completion.
\end{enumerate}
\end{proposition}
\begin{proof}
For (1), the opens \(U=\mathbf A^1_k\setminus D\), with \(D\) finite,
are cofinal.  Their coordinate rings are localizations of the PID
\(k[t]\), so \(\operatorname{Pic}(U)=0\).  Kummer theory gives
\[
 H^1_{\mathrm{cont}}(U,\Z_3(1))
   =\bigoplus_{a\in D}\Z_3,
\]
with coordinates the valuations at \(t-a\).  Here \(k^\times\) is
\(3\)-divisible, and the transition maps on
\(H^0(U,\mu_{3^r})\) are surjective.  Consequently the
\(\varprojlim{}^1\) term in continuous cohomology vanishes.  These are
continuous coefficient systems, not discrete constant
\(\Z_3\)-sheaves
\cite[(3.1), Theorem~3.2 and (3.27), pp.~216--217, 222]{BJannsen1988}.
On the field, \(L^\times/k^\times\) is freely generated by the linear
polynomials \(t-a\).  Therefore the corresponding calculation gives
\[
\begin{split}
 H^1_{\mathrm{cont}}(\Spec L,\Z_3(1))
 &=\varprojlim_r L^\times/(L^\times)^{3^r}\\
 &=\left\{(v_a)\in\prod_{a\in k}\Z_3:
       \text{for every }r,\ v_a\in3^r\Z_3
       \text{ for all but finitely many }a\right\}.
\end{split}
\]
The open colimit is the finite-support subgroup
\(\bigoplus_{a\in k}\Z_3\), proving injectivity.

Enumerate the monic irreducible polynomials \(P_n\in\F_2[t]\),
\(n\geq1\), by degree and then by binary coefficient encoding.  Define
\[
 c_r=\prod_{1\leq n<r}P_n(t)^{3^n}
       \quad\text{in }L^\times/(L^\times)^{3^r}.
\]
Since \(c_{r+1}/c_r=P_r(t)^{3^r}\), these classes are compatible.
Their limit has valuation \(3^n\ne0\) at each geometric root of
\(P_n\), so it has infinite support.  Each \(c_r\) belongs to
\(\F_2(t)^\times\); the resulting Kummer class is Galois invariant.
No nonzero scalar in \(\Q_3\) makes its support finite.  This proves
all assertions in (1), including the rational one.

For (2), take the group of locally constant integer-valued functions
\[
 M=C_{\mathrm{lc}}(\Z_3,\Z),\qquad (Fh)(x)=h(x+1),
\]
where \(\Z\) is discrete.  Each function factors through some
\(\Z/3^s\), and a finite sum of rational multiples of such functions
does too.  Thus \(F^{3^s}h=h\) for some \(s\), and
\(Fh=4h\) in \(M\otimes\Q\) implies
\((4^{3^s}-1)h=0\), hence \(h=0\).

Choosing integer lifts on the finitely many clopen fibres gives
\[
 M/3^rM=C_{\mathrm{lc}}(\Z_3,\Z/3^r),\qquad
 \widehat M_3=C(\Z_3,\Z_3).
\]
The second identity follows by taking compatible reductions and is
\(F\)-equivariant.  The elementary identity
\[
 v_3(4^{3^s}-1)=s+1
\]
follows by induction: if \(u=1+3^a b\), \(a\geq1\) and \(3\nmid b\),
then
\[
 u^3-1=3^{a+1}b\bigl(1+3^a b+3^{2a-1}b^2\bigr),
\]
and the last factor is a \(3\)-adic unit.  Consequently the functions
\[
 h_r(x)=4^{\,x\bmod 3^{r-1}}\pmod{3^r}
\]
are compatible and satisfy \(Fh_r=4h_r\), including at the period's
wraparound.  Their limit \(h\in\widehat M_3\) satisfies
\(Fh=4h\) and \(h(0)=1\).  It is not \(3\)-primary torsion and
remains a nonzero eigenvector after tensoring with \(\Q_3\).
\end{proof}

The second example uses inflation maps between finite permutation
lattices.  It does not invalidate the following distinct calculation.
If \(P=\bigoplus_O\Z[O]\) for fixed finite \(F\)-orbits \(O\),
its \(3\)-adic completion embeds in \(\prod_O\Z_3[O]\).
For a \(q\)-eigenvector, projection to an orbit of size \(d\) gives
\((q^d-1)v_O=0\).  If \(q\in\Q_3^\times\) is not a root of unity,
all projections vanish, even after rationalization.  Thus a residue
argument must examine the actual transition maps and the unramified
kernel; the fixed-orbit permutation calculation itself remains valid.

Neither example identifies its completed group with the degree-four
birational cohomology in B1.  They show precisely why finite-coefficient
norm-residue comparison, or finite Frobenius orbits of individual
symbols, cannot supply the missing interchange without uniform
ramification control.  The exact finite programs are
\begin{quote}
\small
\path{part_b/code/kummer_limit_counterexample.py}\\
\path{part_b/code/frobenius_completion_counterexample.py}
\end{quote}
with their correspondingly named \texttt{.txt} outputs in
\texttt{part\_b/output}.  They check the stated finite congruences;
the infinite assertions are proved above.
